%-----------------------------QUESTÃO------------------------------------
%Cód.: TF3p24q06


\question [\questaopt] Determine a intensidade do campo elétrico criado por
uma carga pontual $ Q $ de \SI{-8,0}{\micro\coulomb}, em um ponto A situado a \SI{6,0}{\micro\coulomb}
dessa carga. \\
\textbf{Dados:} Constante eletrostática do vácuo $ \cteK $

%------------------------------ALTERNATIVAS-----------------------------

\begin{EnvUplevel}
	\begin{choices}
		
		\choice \SI[per-mode = symbol-or-fraction]{6,0e5}{\newton\per\coulomb}  
		\choice \SI[per-mode = symbol-or-fraction]{2,0e4}{\newton\per\coulomb}
		\choice \SI[per-mode = symbol-or-fraction]{4,0e7}{\newton\per\coulomb}
		\correctchoice \SI[per-mode = symbol-or-fraction]{2,0e7}{\newton\per\coulomb}
		\choice \SI[per-mode = symbol-or-fraction]{7,0e7}{\newton\per\coulomb}
		
	\end{choices}
\end{EnvUplevel}

%--------------------------SOLUÇÃO---------------------------------------

\begin{EnvUplevel}
	\begin{solution}[2cm]
		{\footnotesize\\
			A intensidade do campo elétrico criado por uma partícula eletrizada
			é determinada pela relação:
			$$
			\Eel \\
			$$
			Para o ponto A, temos $\SI{6,0}{\centi \meter} = \SI{6,0e-2}{\meter}$.
			Assim:
			$$
			E=\num{9e9}\times\frac{\left| \num{-8,0e-6} \right|}{\left( \num{6,0e-2} \right)^{2}}
			$$
			$$
			F_e=\SI[per-mode = symbol-or-fraction]{2,0e7}{\newton\per\coulomb}
			$$
		}
		
	\end{solution}
\end{EnvUplevel}
